% 张三慧第4版V2第19章《磁场中的磁介质》独立物理解答。
% 原扫描PDF物理页187–189。图读数仅按教材图标近似，专家复核待办。
\section*{第19章\quad 磁场中的磁介质：独立解析}

\paragraph{习题 19.1}
令 \(x=mB/(k_BT)\)。能量
\(W_\uparrow=-mB\)、\(W_\downarrow=+mB\)，
单原子配分函数 \(Z=e^x+e^{-x}\)。
取向两态概率为
\(P_\uparrow=e^x/Z\)、
\(P_\downarrow=e^{-x}/Z\)，
每原子沿磁场方向的平均磁矩
\[
 \langle m_z\rangle=m(P_\uparrow-P_\downarrow)
 =m\tanh x.
\]
单位体积 \(N\) 个原子，所以
\[
 \boxed{M=Nm\tanh\!\left(\frac{mB}{k_BT}\right)
 =Nm\frac{e^{x}-e^{-x}}{e^{x}+e^{-x}}}.
\]
若 \(x\ll1\)，\(\tanh x=x+O(x^3)\)，则
\[
 \boxed{M\simeq\frac{Nm^2B}{k_BT}\propto B/T}.
\]
若 \(x\gg1\)，\(\tanh x\to1\)，有
\(\boxed{M\to Nm}\)（饱和）。
\textbf{源模型边界：}
这是只有两个方向的理想独立取向模型，
不包含真实铁磁交换作用和饱和场自洽修正。

\paragraph{习题 19.2}
轨道电子的磁矩与角动量大小关系
\(\mu=eL/(2m_e)\)，由于电子电量为负，
\(\boldsymbol\mu=-e\boldsymbol L/(2m_e)\)。
外场力矩
\(\boldsymbol\tau=\boldsymbol\mu\times\boldsymbol B\)，
大小
\[
 \boxed{\tau=\frac{eLB}{2m_e}\sin\theta}.
\]
由 \(d\boldsymbol L/dt=\boldsymbol\tau\)，
若 \(\boldsymbol L\) 绕 \(\boldsymbol B\)
以角速率 \(\Omega\) 进动，
\(|d\boldsymbol L/dt|=\Omega L\sin\theta\)，
故
\[
 \boxed{\Omega=\frac{eB}{2m_e}
 \approx8.79\times10^{10}\,\mathrm{rad/s}
 \quad (B=1\,\mathrm T)}.
\]
进动旋向依据带负电荷与图19.2定义判断，
该值为角速度大小。

\paragraph{习题 19.3}
玻尔氢原子模型中
\(r=.53\times10^{-10}\,\mathrm m\)、
\(v=2.2\times10^6\,\mathrm{m/s}\)。
运动电子构成的圆电流大小
\(I=e/T=ev/(2\pi r)\)，
在圆心的磁场
\[
 \boxed{B=\frac{\mu_0I}{2r}
 =\frac{\mu_0ev}{4\pi r^2}
 \approx12.55\,\mathrm T}.
\]
以题给自旋角动量
\(S=0.53\times10^{-34}\,\mathrm{J\,s}\)
及题给磁矩
\(\mu_s=1.41\times10^{-26}\,\mathrm{A\,m^2}\)，
磁场力矩导致进动
\[
 \boxed{\Omega_s=\frac{\mu_s B}{S}
 \approx3.34\times10^9\,\mathrm{rad/s}}.
\]
本模型属于半经典轨道及自旋估算，
真实氢原子结构需量子力学。

\paragraph{习题 19.4}
铁棒半径 \(r=.005\,\mathrm m\)、长
\(L=.12\,\mathrm m\)，体积
\[
 V=\pi r^2L=9.4248\times10^{-6}\,\mathrm{m^3}.
\]
铁质量 \(m_{\rm Fe}=7800V\)，
原子数
\(N_{\rm Fe}=m_{\rm Fe}N_A/.05585\)。
每原子两电子自旋参与且全部平行，
电子数
\[
 \boxed{N_s=2N_{\rm Fe}
 \approx1.59\times10^{24}}.
\]
以一电子一个玻尔磁子
\(\mu_B=9.274\times10^{-24}\,\mathrm{A\,m^2}\)，
总磁矩
\[
 \boxed{\mu_{\rm tot}=N_s\mu_B
 \approx14.7\,\mathrm{A\,m^2}}.
\]
均匀磁化强度 \(M=\mu_{\rm tot}/V
 \approx1.56\times10^6\,\mathrm{A/m}\)。
长圆柱体等效环向面电流密度大小
\(K_m=M\)，
沿长度 \(L\) 的等效总面电流
\[
 \boxed{I_{\rm eq}=ML\approx1.87\times10^5\,\mathrm A}.
\]
理想长螺线管模型的内磁场
\[
 \boxed{B'=\mu_0M\approx1.96\,\mathrm T}.
\]
对有限长棒，上式是忽略端部退磁的估计，
不能直接断言其端面附近场处处同值。

\paragraph{习题 19.5}
磁针体积
\(V=(.085)(.010)(.00020)
 =1.70\times10^{-7}\,\mathrm{m^3}\)。
铁密度及摩尔质量同题，
全部两电子自旋取向一致的
总磁矩
\[
 \boxed{\mu=2\frac{7800V}{.05585}N_A\mu_B
 \approx0.265\,\mathrm{A\,m^2}}.
\]
磁针磁矩与地磁场正交时，
\[
 \boxed{\tau_{\max}=\mu B_E
 \approx1.38\times10^{-5}\,\mathrm{N\,m}}.
\]
磁矩与外场\emph{同向平行}时的势能
\[
 \boxed{U=-\mu B_E
 \approx-1.38\times10^{-5}\,\mathrm J}.
\]
若反平行，势能变为正同样的绝对值。

\paragraph{习题 19.6}
均匀理想螺线环的安培环路给
\[
 \boxed{H=\frac{NI}l=\frac{20(.1)}{.10}
 =20\,\mathrm{A/m}}.
\]
没有介质时
\[
 \boxed{B_0=\mu_0H
 \approx2.51\times10^{-5}\,\mathrm T}.
\]
加入线性相对磁导率
\(\mu_r=4200\) 的介质，
\[
 \boxed{H=20\,\mathrm{A/m}},\quad
 \boxed{B=\mu_0\mu_rH
 \approx1.056\times10^{-1}\,\mathrm T}.
\]
自由电流分量仍
\(B_0=\mu_0H\)，
磁化（束缚）电流贡献
\[
 \boxed{B'=B-B_0
 =\mu_0(\mu_r-1)H
 \approx0.1055\,\mathrm T}.
\]
对真实非线性铁磁体，常数 \(\mu_r\)
只适用于给定工作点的近似。

\paragraph{习题 19.7}
\(S=1\,\mathrm{cm^2}=10^{-4}\,\mathrm{m^2}\)，
磁通 \(\Phi=2\times10^{-6}\,\mathrm{Wb}\)，故
\[
 \boxed{B=\Phi/S=0.0200\,\mathrm T}.
\]
\[
 \boxed{H=\frac{NI}{l}
 =\frac{300(.032)}{.30}=32.0\,\mathrm{A/m}}.
\]
总磁导率与相对磁导率
\[
 \boxed{\mu=\frac BH
 =6.25\times10^{-4}\,\mathrm{H/m}},
\quad
 \boxed{\mu_r=\frac B{\mu_0H}
 \approx497}.
\]
在 \(B=\mu_0(H+M)\) 下
\[
 \boxed{M=\frac B{\mu_0}-H
 \approx1.588\times10^4\,\mathrm{A/m}}.
\]
均匀环向磁化的等效面电流密度
\(\boxed{K_m=M}\)，
单位为 A/m；不是直接把 \(M\) 写成安培。

\paragraph{习题 19.8}
环形线圈1000匝，平均半径 \(r=.15\,\mathrm m\)、
电流 \(2.0\,\mathrm A\)：
\[
 \boxed{H=\frac{NI}{2\pi r}
 =\frac{2000}{.30\pi}
 \approx2.12\times10^3\,\mathrm{A/m}}.
\]
\[
 \boxed{\mu=\frac BH
 \approx4.71\times10^{-4}\,\mathrm{H/m}},\quad
 \boxed{\mu_r=\frac B{\mu_0H}=375}.
\]
材料磁化强度
\[
 \boxed{M=\frac B{\mu_0}-H
 \approx7.94\times10^5\,\mathrm{A/m}}.
\]
磁化面电流密度的大小
\(\boxed{K_m=M}\)。

\paragraph{习题 19.9}
图19.25的纵轴为
\(B/\mathrm T\)，横轴为
\(H/(\mathrm{A\,m^{-1}})\)。
按图估读 \(B=1.2\,\mathrm T\)
时铁芯原来的
\(H_{\rm core}\) 约
\(1.7\times10^2\,\mathrm{A/m}\)
（因线宽和曲线陡升，误差需保留）。
长螺线管匝密度 \(n\) 不变，
原有 \(nI=H_{\rm core}\)；
抽去铁芯后保持相同磁场
\(B=\mu_0nI_{\rm air}\)，所以
\[
 \boxed{\frac{I_{\rm air}}{I_{\rm core}}
 =\frac{B}{\mu_0H_{\rm core}}
 \approx5.6\times10^3}.
\]
若原来 \(I_{\rm core}=6.0\,\mathrm A\)，
\[
 \boxed{I_{\rm air}\approx3.37\times10^4\,\mathrm A}.
\]
\textbf{图读风险：}需按辅导书对应题号
检查 \(H_{\rm core}\) 的准确读数；
这是依图19.25的近似量，不可视为铁芯
始终有固定相对磁导率。

\paragraph{习题 19.10}
长螺线管每米 \(n=800\) 匝，
目标场强 \(B=1.0\,\mathrm T\)。
由教材图19.25初始曲线约读
\(H\sim1.5\times10^2\,\mathrm{A/m}\)。
安培环路给 \(H=nI\)，所以
\[
 \boxed{I\approx\frac{150}{800}
 \approx0.19\,\mathrm A}.
\]
\textbf{精度声明：}
磁化曲线不能仅由
\(B\) 和未知铁芯的常数 \(\mu_r\)
得到 \(H\)；必须从原图读取。
\(0.19\,\mathrm A\) 作为图读近似，
随后要核对配套辅导书的选择与坐标精度。

\paragraph{习题 19.11}
矩形磁滞回线的矫顽力
\(H_c=2\,\mathrm{A/m}\)，
为令磁化翻转，脉冲环向励磁场
至少达到反向矫顽强度（临界意义）。
环外径 \(d_o=.8\,\mathrm{mm}\)、
内径 \(d_i=.5\,\mathrm{mm}\)，
平均直径
\(d_m=(d_o+d_i)/2=.65\,\mathrm{mm}\)，
平均磁路长度
\(l_m=\pi d_m\)。
由环路定理
\[
 \boxed{NI_{\min}\simeq H_cl_m
 =2\pi(.65\times10^{-3})
 \approx4.08\times10^{-3}\,\mathrm{A\,turn}}.
\]
若题图是一匝穿芯导线 \(N=1\)，
\(\boxed{I_{\rm peak,min}\approx4.08\,\mathrm{mA}}\)。
磁滞翻转恰在矫顽力附近，
真实开关为可靠翻转还需裕量；
方向应与原磁化方向反向。

\paragraph{习题 19.12}
环总平均长 \(l=.61\,\mathrm m\)，
空气隙 \(l_g=.01\,\mathrm m\)，
铁芯磁路 \(l_f=.60\,\mathrm m\)，
\(N=1000,I=1.5\,\mathrm A\)，
气隙 \(B=.18\,\mathrm T\)。
假定截面均匀，铁芯内同样的
磁感应强度大小 \(B\)。
安培环路：
\[
 NI=H_fl_f+H_gl_g
 =\frac B{\mu_0\mu_r}l_f+
 \frac B{\mu_0}l_g.
\]
因此
\[
 \boxed{\mu_r=
 \frac{Bl_f}{\mu_0NI-Bl_g}
 \approx1.27\times10^3}.
\]
分母必须为正，否则这个线性均匀磁路
假设或输入数据不自洽。

\paragraph{习题 19.13}
图19.27的铁芯主磁路长
\(l_f=.500\,\mathrm m\)，
气隙长 \(l_g=.020\,\mathrm m\)，
\(\mu_r=5000\)。
假设气隙内与铁芯中
近似有相同 \(B=3\,\mathrm T\)，
安培环路积分为
\[
 NI=H_fl_f+H_gl_g
 =\frac B{\mu_0}
 \left(\frac{l_f}{\mu_r}+l_g\right).
\]
代入：
\[
 \boxed{NI=\frac3{4\pi\times10^{-7}}
 \left(\frac{.500}{5000}+.020\right)
 \approx4.80\times10^4\,\mathrm{A\,turn}}.
\]
气隙提供主要磁阻。
\(3\,\mathrm T\) 可能超过常规铁芯
实际饱和场，题目仍按给定线性
\(\mu_r\) 理想模型计算。

\paragraph{习题 19.14}
电铃铁芯磁路 \(l_f=.144\,\mathrm m\)，
气隙 \(l_g=.002\,\mathrm m\)，
截面积 \(S=.60\,\mathrm{cm^2}
 =6.0\times10^{-5}\,\mathrm{m^2}\)，
磁通 \(\Phi=4.8\times10^{-6}\,\mathrm{Wb}\)，
\[
 \boxed{B=\Phi/S=.0800\,\mathrm T}.
\]
安培环路给
\[
 \boxed{NI=\frac B{\mu_0}
 \left(l_g+\frac{l_f}{\mu_r}\right)
 =\frac{.08}{4\pi\times10^{-7}}
 \left(.002+\frac{.144}{1600}\right)
 \approx1.33\times10^2\,\mathrm{A\,turn}}.
\]
题给电压 \(U=220\,\mathrm V\)、
功率 \(P=20\,\mathrm W\)，
若视作直流稳态线圈输入且忽略铁损，
电流
\(I=P/U\approx.0909\,\mathrm A\)。
故
\[
 \boxed{N=\frac{NI}{I}
 \approx\frac{133.1}{.0909}
 \approx1.46\times10^3\text{ 匝}}.
\]
若电铃用交流电，
功率因数和电感电流不能由
\(P/U\) 唯一确定；
源题采用的定常平均电流模型需明确。
